\documentclass[twoside]{article}

\usepackage[preprint]{aistats2027}
\usepackage[round,authoryear]{natbib}
\usepackage{amsmath,amssymb,amsthm}
\usepackage{mathtools}
\usepackage{microtype}
\usepackage{balance}
\usepackage{url}

\renewcommand{\bibname}{References}
\renewcommand{\bibsection}{\subsubsection*{\bibname}}
\bibliographystyle{apalike}

\newtheorem{theorem}{Theorem}
\newtheorem{lemma}{Lemma}
\newtheorem{corollary}{Corollary}
\newtheorem{proposition}{Proposition}
\newtheorem{definition}{Definition}
\newcommand{\E}{\mathbb E}
\newcommand{\calone}{\operatorname{Cal}_1}
\newcommand{\caltwo}{\operatorname{Cal}_2}
\newcommand{\Reg}{\operatorname{Reg}}
\newcommand{\simplex}{\Delta_K}
\newcommand{\1}{\mathbf 1}

\begin{document}
\runningauthor{}

\twocolumn[
\aistatstitle{The Price of Fidelity: Minimax Multiclass Online Recalibration}

\aistatsauthor{Pahan Dewasurendra}

\aistatsaddress{Johns Hopkins University}
]

\begin{abstract}
An online recalibrator observes a multiclass probability forecast, issues a new forecast, and then observes the outcome. It must become calibrated without losing much Brier score relative to the original forecasts. A recent binary result establishes exponent three, up to logarithms, for calibration error $\varepsilon$ and average excess loss $\varepsilon^2$. For $K$ classes, existing algorithms give the upper exponent $K+1$, but no matching lower bound was known. We prove a sharper Pareto lower bound. For every fixed $K$, if average excess Brier loss is at most $\beta$, then an oblivious truthful instance forces expected canonical calibration at least \[ c_K\min\{1,(T\beta^{(K-1)/2})^{-1/2}\}-C_K\sqrt{\beta}. \] Consequently, $(\varepsilon,\varepsilon^2)$ multiclass recalibration needs $\Omega_K(\varepsilon^{-(K+1)})$ rounds. Combined with recent upper bounds, this establishes the fixed-$K$ minimax horizon up to logarithmic factors. The proof introduces a dimension-matched martingale argument. A fractional Burkholder moment converts adaptive bin occupancies into fluctuation, while a $d$-dimensional lattice transport inequality shows that concentrating truthful hints into fewer forecast bins costs squared loss. The result holds for randomized learners and extends to losses that are strongly proper in Euclidean norm.
\end{abstract}

\section{Introduction}

Calibration asks whether empirical outcome frequencies agree with issued probabilities on rounds sharing the same forecast. It can be enforced online against arbitrary outcome sequences, but a calibrated postprocessor can destroy the accuracy of a useful base forecaster. Online recalibration adds a fidelity requirement: after seeing a hint $q_t$, the learner predicts $p_t$ and seeks both calibration and small excess loss relative to $q_t$ \citep{kuleshov2017estimating,okoroafor2024faster}.

The sharp binary picture appeared only recently. \citet{hu2026optimal} prove that expected calibration $\varepsilon$ and average excess squared loss $\varepsilon^2$ are jointly attainable in $\widetilde O(\varepsilon^{-3})$ rounds and require $\Omega(\varepsilon^{-3})$ rounds. Their theory is binary. Their multiclass experiment first projects top-class correctness to a binary event, and they explicitly note that the binary theorem does not apply to the lifted forecast.

For $K\geq3$, the upper rate is nevertheless known. The multiclass construction of \citet{chen2026calibeating} tracks an arbitrary reference predictor while controlling cumulative Brier calibration. With a simplex discretization of scale $\varepsilon$, it yields the horizon $\widetilde O_K(\varepsilon^{-(K+1)})$. The exponent contains the full simplex dimension $d=K-1$. It was unclear whether this dependence was a limitation of the discretized upper bound or an unavoidable cost of preserving the hints.

We show that it is unavoidable. Our main result gives a continuum of lower bounds indexed by the fidelity budget $\beta$. At the canonical choice $\beta=\varepsilon^2$, it proves the matching exponent $d+2=K+1$. The hard hints are not inaccurate or adversarially corrupted. They are the true conditional outcome probabilities. Properness then turns excess expected loss into squared movement away from the hints. A learner can merge nearby hints to reduce the number of calibration bins, but doing so spends this movement budget. The theorem quantifies the resulting geometric tradeoff.

The main technical obstacle is adaptivity. A bin is selected using previous outcomes, so its occupancy and its label fluctuation are dependent. Standard independent-sum estimates do not apply after conditioning on the final bin. We combine two pathwise and martingale ingredients:
\begin{itemize}
\item A fractional Burkholder--Rosenthal argument at moment
$2(1+2/d)$ lower bounds total fluctuation using the $(1+2/d)$-moment of adaptive occupancies.
\item A lattice transport inequality upper bounds exactly that occupancy
moment by the squared distance used to merge a $d$-dimensional grid of truthful hints.
\end{itemize}
The moment is chosen to match the geometry. The identities behind the final rate are exact, not an optimization artifact.

Our result concerns canonical calibration of pure forecasts for fixed $K$ as $\varepsilon$ vanishes. This differs from recent high-dimensional calibration work, which studies calibration without reference fidelity and optimizes the dependence on $K$ for fixed accuracy \citep{peng2025highdim,fishelson2025highdim}. It also differs from U-calibration, projected smooth calibration, and RKHS multicalibration, which weaken or change the exact-bin canonical criterion \citep{gopalan2024computational,luo2024optimal,farina2026defensive}.

\section{Setting}

Let $K\geq2$, let $d=K-1$, and let $\simplex=\{p\in\mathbb R_+^K:\sum_i p_i=1\}$. On round $t$, an oblivious environment provides a hint $q_t\in\simplex$. The learner observes $q_t$ and outputs $p_t\in\simplex$, possibly using randomization and the past. It then observes $y_t\in\{e_1,\ldots,e_K\}$. We use Brier loss $\ell(p,y)=\|p-y\|_2^2$ \citep{brier1950verification} and normalized excess loss \[ \Reg_T(p,q,y)=\frac1T\sum_{t=1}^T \bigl(\ell(p_t,y_t)-\ell(q_t,y_t)\bigr). \]

For a realized prediction sequence, let $V(p)$ be its distinct values. Write $n_v=\sum_t\1\{p_t=v\}$ and $\rho_v=n_v^{-1}\sum_{t:p_t=v}y_t$ for occupied $v$. We study canonical mean and root-mean-square calibration
\begin{align*}
 \calone(p,y)&=\sum_{v\in V(p)}\frac{n_v}{T}\|v-\rho_v\|_2,\\
 \caltwo(p,y)^2&=\sum_{v\in V(p)}\frac{n_v}{T}\|v-\rho_v\|_2^2.
\end{align*}
Weighted Cauchy--Schwarz gives $\calone\leq\caltwo$. Moreover, $T\caltwo^2$ is exactly the cumulative Brier calibration used by \citet{chen2026calibeating}.

An algorithm is $(\alpha,\beta)$-recalibrated at horizon $T$ if, for every deterministic oblivious sequence $(q_t,y_t)_{t\leq T}$, \[ \E\caltwo(p,y)\leq\alpha, \qquad \E\Reg_T(p,q,y)\leq\beta, \] where the expectation is over learner randomization. Using $\calone$ gives a weaker requirement, so a lower bound for $\calone$ also applies to $\caltwo$.

\section{Main Results}

Our first theorem is a Pareto lower bound. Its stochastic instance is a proof device. Corollary~\ref{cor:deterministic} converts it to the standard deterministic-sequence guarantee.

\begin{theorem}[Price of fidelity]\label{thm:main}
Fix $K\geq2$. There are constants $c_K,C_K,\beta_K>0$ such that for every $T\geq1$ and $\beta\in(0,\beta_K)$, there is a deterministic hint sequence $q_{1:T}$ and a product distribution over $y_{1:T}$ with $\E[y_t\mid q_t]=q_t$ for which every nonanticipating randomized learner satisfying $\E\Reg_T\leq\beta$ obeys \[ \E\calone(p,y) \geq c_K\min\!\left\{1, \frac{1}{\sqrt{T\beta^{d/2}}}\right\} -C_K\sqrt\beta. \]
\end{theorem}

The dimension in Theorem~\ref{thm:main} is intrinsic rather than ambient. The supplement proves the following structured form by projecting learner forecasts onto the convex hint box before rounding.

\begin{proposition}[Intrinsic hint dimension]\label{prop:intrinsic}
Fix $K\geq2$ and $1\leq r\leq K-1$. There is an $r$-dimensional affine box $Q_r\subset\operatorname{int}(\simplex)$ such that restricting all truthful hints to $Q_r$ changes the lower bound in Theorem~\ref{thm:main} to \[ \E\calone(p,y) \geq c_{K,r}\min\left\{1, \frac1{\sqrt{T\beta^{r/2}}}\right\}-C_{K,r}\sqrt\beta. \]
\end{proposition}

The constants may depend on fixed $K$. Appendix~A states the corresponding dependence for any loss that is strongly proper in Euclidean norm.

\begin{corollary}[Minimax horizon]\label{cor:deterministic}
For every fixed $K\geq2$ and all sufficiently small $\varepsilon$, any algorithm that is $(\varepsilon,\varepsilon^2)$-recalibrated for all deterministic oblivious sequences must have \[ T=\Omega_K(\varepsilon^{-(K+1)}). \] The conclusion holds with either $\calone$ or $\caltwo$ calibration.
\end{corollary}

\begin{proof}
Average a claimed deterministic-sequence guarantee over the product distribution in Theorem~\ref{thm:main}. If both expected guarantees held, the theorem with $\beta=\varepsilon^2$ would imply \[ \varepsilon\geq c_K\min\{1,(T\varepsilon^d)^{-1/2}\}-C_K\varepsilon. \] For small enough $\varepsilon$, this requires $T\geq c'_K\varepsilon^{-(d+2)}$. Since $d+2=K+1$, the claim follows. The inequality $\calone\leq\caltwo$ gives the RMS statement.
\end{proof}

For completeness, the known upper bound can be expressed in our notation.

\begin{proposition}[Existing upper bound]\label{prop:upper}
For fixed $K\geq3$, there is an algorithm such that for every $\eta\in(0,1)$ and every oblivious sequence,
\begin{align*}
 \E\Reg_T&=O_K(\eta^2),\\
 \E\caltwo&=\widetilde O_K\!\left(
 T^{-1/4}+\eta+\frac1{\sqrt{T\eta^d}}\right).
\end{align*}
Thus $(\varepsilon,\varepsilon^2)$-recalibration is attainable when $T=\widetilde O_K(\varepsilon^{-(K+1)})$. For $K=2$, the same horizon $\widetilde O(\varepsilon^{-3})$ follows from \citet{hu2026optimal}.
\end{proposition}

\begin{proof}
Theorem 17 of \citet{chen2026calibeating}, with the reference algorithm set to $q_t$, gives expected cumulative regret $O(\eta^2T)$ and cumulative Brier calibration $\widetilde O_K(\sqrt T+\eta^{-d}+\eta^2T)$ with high probability. Integrating the tail and applying Jensen gives the display. Set $\eta=\varepsilon$. When $d\geq2$ and $T=\widetilde\Omega_K(\varepsilon^{-(d+2)})$, all three calibration terms are $\widetilde O_K(\varepsilon)$.
\end{proof}

Together, Corollary~\ref{cor:deterministic} and Proposition~\ref{prop:upper} establish the fixed-$K$ horizon up to logarithms. The sharper lower bound also matches the dimension-dependent term in the upper tradeoff after identifying $\eta=\sqrt\beta$.

Define $\varepsilon_K^*(T)$ as the infimum over $\varepsilon$ for which some algorithm is $(\varepsilon,\varepsilon^2)$-recalibrated at horizon $T$. Equivalently, this is the smallest joint scale when calibration is measured in units of $\varepsilon$ and fidelity in units of $\varepsilon^2$.

\begin{corollary}[Joint minimax rate]\label{cor:rate}
For every fixed $K\geq2$, \[ \varepsilon_K^*(T)=\widetilde\Theta_K \bigl(T^{-1/(K+1)}\bigr). \] The lower bound has no logarithmic loss.
\end{corollary}

\begin{proof}
Invert the horizon bounds in Corollary~\ref{cor:deterministic} and Proposition~\ref{prop:upper}.
\end{proof}

\section{Truthful Lattice Instance}

We now prove Theorem~\ref{thm:main}. Choose constants $0<a_K<b_K$ so that \[ Q=\{q(x)=(x_1,\ldots,x_d,1-\textstyle\sum_{j=1}^d x_j): x\in[a_K,b_K]^d\} \] lies inside $\simplex$. We may take $a_K=1/(4K)$ and $b_K=1/(2K)$. In particular, $q_1(1-q_1)\geq\sigma_K^2>0$ throughout $Q$.

Let $G_h$ be the image of the spacing-$h$ lattice in this box. Its size $M$ satisfies $c_Kh^{-d}\leq M\leq C_Kh^{-d}$. Cycle deterministically through $G_h$. Each source point then occurs at most $L=\lceil T/M\rceil$ times. Draw labels independently as \[ y_t\mid q_t\sim\operatorname{Categorical}(q_t). \] Represent learner randomization by a seed $U$ sampled at time zero, independently of the labels, and set $\mathcal F_t=\sigma(U,y_1,\ldots,y_t)$. Then $p_t$ is $\mathcal F_{t-1}$-measurable and $\E[y_t\mid\mathcal F_{t-1}]=q_t$. Truthfulness gives the exact identity \[ \E[\ell(p_t,y_t)-\ell(q_t,y_t)\mid\mathcal F_{t-1}] =\|p_t-q_t\|_2^2. \] Hence the fidelity premise implies
\begin{equation}\label{eq:movement-p}
 \E\sum_{t=1}^T\|p_t-q_t\|_2^2\leq\beta T.
\end{equation}

\section{Rounding and Fluctuation}

We first reduce arbitrary forecasts to a finite net. Let $m=\lceil\sqrt K/h\rceil$ and \[ N_h=\{z/m:z\in\mathbb Z_+^K,\ \sum_i z_i=m\}. \] Largest-remainder rounding defines $R_h:\simplex\to N_h$ with $\|R_h(p)-p\|_2\leq h$ and $|N_h|\leq C_Kh^{-d}$. Put $r_t=R_h(p_t)$.

\begin{lemma}[Rounding]\label{lem:rounding}
Pathwise, $\calone(r,y)\leq\calone(p,y)+h$. Moreover, \[ \E B:=\E\sum_{t=1}^T\|r_t-q_t\|_2^2 \leq C_K(\beta+h^2)T. \]
\end{lemma}

\begin{proof}
For each $v\in N_h$, group all exact values $p$ rounded to $v$. Then
\begin{align*}
 \left\|\sum_{t:r_t=v}(v-y_t)\right\|_2
 &\leq\sum_{p:R_h(p)=v}
 \left\|\sum_{t:p_t=p}(p-y_t)\right\|_2\\
 &\quad+\sum_{t:r_t=v}\|v-p_t\|_2.
\end{align*}
Sum over $v$ and divide by $T$. The movement claim follows from $\|r_t-q_t\|_2^2\leq2\|p_t-q_t\|_2^2+2h^2$ and \eqref{eq:movement-p}.
\end{proof}

Set $h=c\sqrt\beta$ for a sufficiently small fixed $c$. Constants below absorb $c$. Lemma~\ref{lem:rounding} gives $\E B\leq C_Kh^2T$. For every $v\in N_h$, define \[ \nu_v=\sum_t\1\{r_t=v\}, \qquad \mu_v=\sum_t\1\{r_t=v\}(q_{t,1}-y_{t,1}). \]

\begin{lemma}[Bias and fluctuation]\label{lem:bias}
\[ \E\calone(r,y)\geq \frac1T\sum_v\E|\mu_v|-\sqrt{\E B/T}. \]
\end{lemma}

\begin{proof}
The first coordinate of the unnormalized residual in bin $v$ equals \[ \mu_v+\sum_{t:r_t=v}(r_{t,1}-q_{t,1}). \] The Euclidean norm dominates its absolute value. The reverse triangle inequality leaves $\sum_v|\mu_v|$ minus the total absolute bias. Two Cauchy--Schwarz inequalities bound the latter pathwise by $\sqrt{TB}$. Take expectations and apply Jensen.
\end{proof}

\section{A Dimension-Matched Moment}

Set \[ s=1+\frac2d, \qquad a=\frac{2s-1}{2s-2}, \qquad b=\frac1{2s-2}. \] The identities
\begin{equation}\label{eq:identities}
 a-b=1,
 \qquad a-sb=\frac12,
 \qquad 2b=\frac d2
\end{equation}
will determine the final exponent.

\begin{lemma}[Adaptive anti-concentration]\label{lem:anti}
For the predictable bin choices above, \[ \sum_v\E|\mu_v| \geq c_K\frac{T^a}{(\sum_v\E\nu_v^s)^b}. \]
\end{lemma}

\begin{proof}
Fix $v$ and write $X_t=\1\{r_t=v\}(q_{t,1}-y_{t,1})$. These are bounded martingale differences. Orthogonality and the interior variance bound give
\begin{equation}\label{eq:second}
 \E\mu_v^2\geq\sigma_K^2\E\nu_v.
\end{equation}
Their predictable quadratic variation is at most $\nu_v$. The Burkholder--Rosenthal inequality at the real moment $2s\geq2$ gives
\begin{align*}
 \E|\mu_v|^{2s}
 &\leq C_s\E\left[
 \langle\mu_v\rangle_T^s+\max_t|X_t|^{2s}\right]\\
 &\leq C_K\E\nu_v^s.
\end{align*}
The second term is absorbed because it is at most $\Pr(\nu_v>0)\leq\E\nu_v^s$. Interpolating $L_2$ between $L_1$ and $L_{2s}$, then using \eqref{eq:second}, yields \[ \E|\mu_v|\geq c_K\frac{(\E\nu_v)^a}{(\E\nu_v^s)^b}. \] Omit bins with zero expected occupancy. Since $a-1=b$, H\"older's inequality gives \[ T=\sum_v\E\nu_v \leq\left(\sum_v\frac{(\E\nu_v)^a}{(\E\nu_v^s)^b}\right)^{1/a} \left(\sum_v\E\nu_v^s\right)^{b/a}. \] Rearranging proves the claim. The noninteger moment is valid in the Burkholder--Rosenthal theorem for every real exponent at least two \citep{burkholder1973distribution,collina2026optimal}.
\end{proof}

The use of $s=1+2/d$ is essential. A fourth-moment argument corresponds to $s=2$ and aligns directly with squared $d$-dimensional transport only when $d=2$. The binary proof of \citet{hu2026optimal} instead uses a separate one-dimensional charging argument.

\section{Lattice Transport}

For an output bin $v$, define its movement contribution $b_v=\sum_{t:r_t=v}\|q_t-v\|_2^2$.

\begin{lemma}[Occupancy transport]\label{lem:transport}
Pathwise, for every $v\in N_h$, \[ \nu_v^{1+2/d} \leq C_K\left[ \1\{\nu_v>0\}L^{1+2/d} +\frac{L^{2/d}b_v}{h^2}\right]. \]
\end{lemma}

\begin{proof}
Sort source occurrences assigned to $v$ by distance. A radius-$R$ Euclidean ball centered at any $v$, including a boundary point, contains at most $C_KL(1+R/h)^d$ occurrences of the spacing-$h$ source lattice. After the first $C_KL$ assignments, the $i$th distance is therefore at least $c_Kh(i/L)^{1/d}$. Summing the squared ordered distances and comparing with an integral gives \[ b_v\geq c_Kh^2L^{-2/d} \bigl(\nu_v^{1+2/d}-C_KL^{1+2/d}\bigr). \] Rearrange.
\end{proof}

\begin{lemma}[Global occupancy moment]\label{lem:global}
If $\E B\leq C_Kh^2T$, then \[ \sum_v\E\nu_v^s\leq
 \begin{cases}
 C_KT,&T<M,\\
 C_KT^sh^2,&T\geq M.
 \end{cases}
\]
\end{lemma}

\begin{proof}
Sum Lemma~\ref{lem:transport} over occupied output bins. If $T<M$, then $L=1$ and there are at most $T$ such bins. The baseline is $O_K(T)$ and the expected transport term is at most $\E B/h^2=O_K(T)$.

If $T\geq M$, then $L\leq C_KTh^d$ and at most $|N_h|\leq C_Kh^{-d}$ bins are occupied. The baseline is at most \[ C_Kh^{-d}(Th^d)^s=C_KT^sh^2, \] because $d(s-1)=2$. The expected transport term is at most \[ C_K(Th^d)^{2/d}\frac{\E B}{h^2} \leq C_KT^{1+2/d}h^2=C_KT^sh^2. \]
\end{proof}

\section{Completing the Lower Bound}

When $T<M$, Lemmas~\ref{lem:anti} and \ref{lem:global} give $\sum_v\E|\mu_v|\geq c_KT$. When $T\geq M$, they give
\begin{align*}
 \sum_v\E|\mu_v|
 &\geq c_K\frac{T^a}{(T^sh^2)^b}\\
 &=c_K\sqrt T\,h^{-d/2},
\end{align*}
where the last equality uses \eqref{eq:identities}. Lemmas~\ref{lem:rounding} and \ref{lem:bias} now imply \[ \E\calone(p,y) \geq c_K\min\{1,(Th^d)^{-1/2}\}-C_Kh. \] Substituting $h=c\sqrt\beta$ proves Theorem~\ref{thm:main}.

The proof is uniform over randomized learners. The seed construction makes every adaptive bin selection predictable. Finally, if an algorithm met both guarantees for every deterministic sequence, averaging over the product label distribution would preserve both guarantees. Therefore some deterministic realization must violate one of them, which justifies the oblivious minimax form in Corollary~\ref{cor:deterministic}.

\section{Related Work and Discussion}

Classical online calibration began with binary forecasting \citep{foster1998asymptotic}. Recalibration adds comparison to a supplied forecaster. \citet{kuleshov2017estimating} introduced this online objective, and \citet{okoroafor2024faster} used approachability to improve its binary calibration--regret tradeoff. \citet{hu2026optimal} obtained the sharp binary $(\varepsilon,\varepsilon^2)$ law and its lower bound. Our construction starts from the same truthful-hint principle, but the proof must control adaptive merging over a full simplex. The fractional moment and lattice transport lemmas provide that step.

Calibeating compares against the refinement of external forecasts rather than their realized loss \citep{foster2023calibeating}. Recent work reduces it to online regret minimization and extends it across proper losses \citep{chen2026calibeating,foster2026proper,fichtl2026bregman}. Theorem 17 of \citet{chen2026calibeating} is particularly relevant because its meta-algorithm also tracks the realized loss of an arbitrary reference predictor. It supplies the matching multiclass upper bound in Proposition~\ref{prop:upper}. Their work does not give a fidelity-constrained multiclass lower bound.

\citet{fishelson2025full} develop the full-swap and triangulation machinery behind dimension-dependent calibration upper bounds. Their lossless rounding analysis is one-dimensional, and they leave its extension to higher dimensions open for general convex losses. Our result is a lower bound rather than a rounding algorithm. It shows that the simplex discretization exponent in the Brier fidelity tradeoff cannot be removed.

High-dimensional calibration has recently taken a different route. \citet{peng2025highdim} and \citet{fishelson2025highdim} obtain calibration-only algorithms with polynomial dependence on the number of classes for fixed accuracy, together with lower bounds when dimension grows with accuracy. These results use no external-hint fidelity constraint and address a different asymptotic regime. U-calibration achieves $\Theta(\sqrt{KT})$ regret uniformly over bounded proper losses without requiring canonical calibration \citep{luo2024optimal}. Computational multiclass work similarly studies weaker projected or smooth criteria \citep{gopalan2024computational}. None implies the fixed-$K$, sharp-in-$\varepsilon$ law proved here.

Our theorem leaves two natural directions. First, the constants are not optimized in $K$, so the result should not be read as a sharp high-dimensional bound. Proposition~\ref{prop:intrinsic} suggests adapting upper bounds to the intrinsic dimension of structured hint sequences. Second, the upper and lower Pareto bounds match in their geometric term, but the upper contains an additional $T^{-1/4}$ pure-sampling term. Determining the complete two-parameter frontier away from $\beta=\Theta(\varepsilon^2)$ remains open.

\section{Conclusion}

Preserving a truthful multiclass forecaster has an unavoidable geometric cost. To reduce the number of exact forecast bins, a recalibrator must transport a $d$-dimensional lattice of hints, and strong propriety charges squared distance for that transport. Adaptive martingale fluctuation then forces calibration error of order $(T\beta^{d/2})^{-1/2}$ before the movement bias is paid. At $\beta=\varepsilon^2$, this gives the minimax horizon $\widetilde\Theta_K(\varepsilon^{-(K+1)})$. The exponent is therefore the price of fidelity, not merely a feature of known simplex discretizations.

\balance
\bibliography{references}

\appendix
\newpage
\section{Pointers to the Supplement}
The supplement gives expanded proofs, the strongly proper loss extension, and the precise conversion of the existing cumulative-calibration upper bound.

\end{document}
